\begin{rubric}{Experiências}
\entry*[2025 -- presente]%
    \textbf{Analista de Tecnologia da Informação,} Dataprev.
    \par Atuação com foco em Arquitetura de Software e DevOps (CI/CD). Responsável pelo desenho e sustentação de soluções para sistemas governamentais de alta criticidade e grande escala, bem como pela implementação, automação e governança de pipelines de integração e entrega contínuas (CI/CD), visando confiabilidade, padronização e resiliência operacional.

\entry*[2020 -- 2022]%
    \textbf{Pesquisador Científico,} Pontifícia Universidade Católica de Minas Gerais (PUC-MG).
    \par Pesquisa aplicada em sistemas operacionais e computação de alto desempenho. Responsável pelo porte do microkernel Nanvix para a arquitetura RISC-V (plataforma PULP) em C e Assembly, resultando em publicação premiada com \href{https://drive.google.com/file/d/1EJsnv8WvVObpRNreRR2fV2hepMs9ifsz/view?usp=sharing}{Menção Honrosa no WSCAD 2022}.

\entry*[2022]%
    \textbf{Monitor de Algoritmos e Estruturas de Dados,} Pontifícia Universidade Católica de Minas Gerais (PUC-MG).
    \par Suporte técnico e acadêmico a alunos de graduação nos tópicos de estruturas de dados (listas, árvores, grafos) e complexidade de algoritmos ($O(n)$), conduzindo plantões de dúvidas e apoio em laboratórios práticos.
\end{rubric}
